% SOURCE_AUTHORITY: sga5_fr_workpass.tex, lines 12621--12635 (LF-normalized unit).
% SOURCE_SHA256: 791F4EFFC5E02832D5D77ED03518C8156D6F07E4C8238B03545DB93D883FBB28
\subsubsection*{4.3}

Observe-se que os módulos assim obtidos não dependem, salvo isomorfismo, senão do ponto $y\in C$, e não da escolha do ponto $y'$ nem de $C'$ acima de $y$. Seja, pois, $y''$ outro ponto dessa natureza. Existe então $g\in G$ tal que $y''=g y'$. A situação $(C',y'',G)$ deduz-se, portanto, da situação $(C',y',G)$ mediante o isomorfismo $g$ sobre $C'$ e $\operatorname{int}(g):h\mapsto ghg^{-1}$ sobre $G$. Conclui-se por transporte de estrutura que os módulos $Sw_y$, etc., associados a $y''$ se deduzem dos associados a $y'$ por extensão dos operadores mediante $\operatorname{int}(g):G\to G$. Sabe-se que essa extensão transforma um módulo em um módulo isomorfo.

\begin{statement}{Proposição 4.4}
Com as notações de 3.7, tomando $(\Lambda=\mathbb Z_\ell)$, e sendo $M^\bullet$ um objeto de $D^+(A[G])$, temos isomorfismos:
\[
(*)\qquad Sw_y\simeq \check{Sw}_y,
\]
\[
(**)\qquad
\Hom^\bullet_{\mathbb Z_\ell[G]}(Sw_y,M^\bullet)
\simeq Sw_y\otimes_{\mathbb Z_\ell[G]}M^\bullet.
\]
\end{statement}
